\documentclass[12pt,a4paper]{article}
\usepackage[utf8]{vietnam}
\usepackage{amsmath,amssymb}
\usepackage{tabularx}
\usepackage{tikz}
\usetikzlibrary{calc,angles,quotes,arrows.meta,patterns}
\usepackage{xcolor}
\usepackage{enumitem}
\usepackage{fancyhdr}
\usepackage{graphicx}
\usepackage[top=0.9cm, bottom=0.9cm, left=1.2cm, right=1.2cm]{geometry}
\usepackage[dethi]{ex_test}
\SAOPTN{widthSA=3.2,heightSA=0.62}

\setlength{\headheight}{15pt}
\setlength{\parskip}{0pt}

\pagestyle{fancy}
\fancyhf{}
\rhead{\textit{Môn Toán 9 -- Đề số 01}}
\lhead{\textit{Bộ đề ôn tập Giữa học kì I}}
\cfoot{\thepage}

\begin{document}

% --- KHUNG TIÊU ĐỀ ĐỀ THI ---
\noindent
\begin{tabularx}{\linewidth}{X r}
\textbf{TRƯỜNG THCS: \dotfill} & \textbf{ĐỀ KIỂM TRA ĐỊNH KỲ GIỮA HỌC KỲ I} \\
\textbf{Họ và tên: \dotfill} & \textbf{NĂM HỌC 2026 -- 2027} \\
\textbf{Lớp: \dotfill \quad SBD: \dotfill} & \textbf{Môn: TOÁN 9} \\
\hline
\end{tabularx}
\vspace{2mm}

\begin{center}
  \textbf{\large ĐỀ ÔN TẬP GIỮA HỌC KỲ I -- ĐỀ SỐ 01}\\
  \textit{Thời gian làm bài: 90 phút (Không kể thời gian phát đề)}
\end{center}
\vspace{1mm}

% --- PHẦN I ---
\noindent\textbf{\large PHẦN I. CÂU TRẮC NGHIỆM NHIỀU PHƯƠNG ÁN LỰA CHỌN (3,0 điểm)}

\textit{Thí sinh trả lời từ câu 1 đến câu 12. Mỗi câu hỏi thí sinh chỉ chọn một phương án duy nhất.}

\vspace{1mm}
\begin{ex}
Phương trình nào sau đây là phương trình bậc nhất hai ẩn?
\choice
{$2x^2 - y = 3$}
{\True $3x + y = 5$}
{$0x + 0y = 4$}
{$\frac{2}{x} + y = 1$}
\end{ex}

\begin{ex}
Cặp số nào sau đây là nghiệm của phương trình $2x - y = 3$?
\choice
{$(1; 1)$}
{\True $(2; 1)$}
{$(0; -3)$}
{$(3; 1)$}
\end{ex}

\begin{ex}
Nghiệm của hệ phương trình $\begin{cases} x + y = 4 \\ x - y = 2 \end{cases}$ là:
\choice
{\True $(3; 1)$}
{$(1; 3)$}
{$(2; 2)$}
{$(4; 0)$}
\end{ex}

\begin{ex}
Điều kiện xác định của phương trình $\frac{x+1}{x-2} = \frac{1}{x+1}$ là:
\choice
{\True $x \neq 2$ và $x \neq -1$}
{$x \neq 2$}
{$x \neq -1$}
{$x \neq 1$ và $x \neq 2$}
\end{ex}

\begin{ex}
Tập nghiệm của bất phương trình $2x - 4 > 0$ là:
\choice
{\True $x > 2$}
{$x < 2$}
{$x \ge 2$}
{$x \le 2$}
\end{ex}

\begin{ex}
Cho bất đẳng thức $-3 < 5$. Khi nhân cả hai vế với $-2$, ta được bất đẳng thức nào sau đây?
\choice
{\True $6 > -10$}
{$6 < -10$}
{$-6 > -10$}
{$-6 < -10$}
\end{ex}

\begin{ex}
Cho tam giác $ABC$ vuông tại $A$, đường cao $AH$. Biết hình chiếu của các cạnh góc vuông trên cạnh huyền là $BH = 4\text{ cm}$ và $HC = 9\text{ cm}$. Độ dài đường cao $AH$ là:
\choice
{\True $6\text{ cm}$}
{$36\text{ cm}$}
{$13\text{ cm}$}
{$5\text{ cm}$}
\end{ex}

\begin{ex}
Cho tam giác $ABC$ vuông tại $A$, đường cao $AH$. Hệ thức lượng nào sau đây đúng?
\choice
{$AB^2 = BH \cdot BC$}
{$AH^2 = BH \cdot CH$}
{$AC^2 = HC \cdot BC$}
{\True Cả ba hệ thức trên đều đúng}
\end{ex}

\begin{ex}
Cho tam giác $ABC$ vuông tại $A$ có độ dài các cạnh góc vuông $AB = 3\text{ cm}$ và $AC = 4\text{ cm}$. Độ dài cạnh huyền $BC$ là:
\choice
{\True $5\text{ cm}$}
{$7\text{ cm}$}
{$25\text{ cm}$}
{$\sqrt{7}\text{ cm}$}
\end{ex}

\begin{ex}
Trong các phương trình sau, phương trình nào là phương trình tích?
\choice
{\True $(x - 1)(x + 2) = 0$}
{$\frac{1}{x} + 2 = 3$}
{$2x + 3 = 0$}
{$x^2 + 1 = 0$}
\end{ex}

\begin{ex}
Cho tam giác $ABC$ vuông tại $A$, có góc $\widehat{B} = 60^\circ$ và cạnh góc vuông $AC = 6\sqrt{3}\text{ cm}$. Độ dài cạnh góc vuông $AB$ là:
\choice
{\True $6\text{ cm}$}
{$12\text{ cm}$}
{$3\text{ cm}$}
{$6\sqrt{3}\text{ cm}$}
\end{ex}

\begin{ex}
Cho hệ phương trình $\begin{cases} 2x + y = 5 \\ x - y = 1 \end{cases}$. Giá trị của tổng $x + y$ là:
\choice
{\True $3$}
{$2$}
{$4$}
{$5$}
\end{ex}


\vspace{1mm}
\noindent\textbf{BẢNG TRẢ LỜI TRẮC NGHIỆM PHẦN I:}
\begin{center}
\begin{tabular}{|c|c|c|c|c|c|c|c|c|c|c|c|c|}
\hline
\textbf{Câu} & \textbf{1} & \textbf{2} & \textbf{3} & \textbf{4} & \textbf{5} & \textbf{6} & \textbf{7} & \textbf{8} & \textbf{9} & \textbf{10} & \textbf{11} & \textbf{12} \\
\hline
\textbf{Chọn} & \rule{0pt}{0.65cm} & & & & & & & & & & & \\
\hline
\end{tabular}
\end{center}
\vspace{2mm}

\setcounter{ex}{0}
\noindent\textbf{\large PHẦN II. CÂU TRẮC NGHIỆM ĐÚNG SAI (2,0 điểm)}

\textit{Thí sinh trả lời câu 1 và câu 2. Trong mỗi ý a), b), c), d) ở mỗi câu, thí sinh chọn đúng hoặc sai.}

\vspace{1mm}
\begin{ex}
Cho hệ phương trình $\begin{cases} 2x - y = 1 \\ x + y = 5. \end{cases}$ Xét tính Đúng/Sai của các khẳng định sau:
\choiceTF
{\True Cặp số $(2; 3)$ là nghiệm của phương trình thứ nhất trong hệ}
{\True Cộng từng vế hai phương trình của hệ ta được phương trình mới $3x = 6$}
{\True Hệ phương trình có nghiệm duy nhất là $(2; 3)$}
{Giá trị của biểu thức $T = x^2 + y^2$ tại nghiệm của hệ bằng $15$}
\end{ex}

\begin{ex}
Cho tam giác $ABC$ vuông tại $A$, đường cao $AH$. Biết $BH = 4\text{ cm}$ và $CH = 9\text{ cm}$. Xét tính Đúng/Sai của các khẳng định sau:
\choiceTF
{\True Độ dài đường cao $AH$ bằng $6\text{ cm}$}
{\True Độ dài cạnh $AB$ bằng $\sqrt{52}\text{ cm}$}
{Hình chiếu của cạnh $AB$ trên cạnh huyền $BC$ bằng $9\text{ cm}$}
{\True Giá trị của $\sin \widehat{ABC}$ bằng $\frac{3}{\sqrt{13}}$}
\end{ex}


\noindent\textbf{BẢNG TRẢ LỜI PHẦN II:}
\begin{center}
\begin{tabular}{|c|c|c|c|c||c|c|c|c|c|}
\hline
\textbf{Câu 1} & \textbf{a)} & \textbf{b)} & \textbf{c)} & \textbf{d)} & \textbf{Câu 2} & \textbf{a)} & \textbf{b)} & \textbf{c)} & \textbf{d)} \\
\hline
\textbf{Đ/S} & \rule{0pt}{0.6cm} & & & & \textbf{Đ/S} & & & & \\
\hline
\end{tabular}
\end{center}
\vspace{2mm}

\setcounter{ex}{0}
\noindent\textbf{\large PHẦN III. CÂU HỎI TRẢ LỜI NGẮN (2,0 điểm)}

\textit{Thí sinh trả lời câu 1 và câu 2. Điền kết quả vào ô trống tương ứng.}

\vspace{1mm}
\begin{ex}
Rút gọn biểu thức $A = \sqrt{12} - 2\sqrt{27} + 3\sqrt{75}$. \shortans{$11\sqrt{3}$}
\end{ex}

\begin{ex}
Cho tam giác $ABC$ vuông tại $A$, đường cao $AH$. Biết $BH = 4\text{ cm}$ và $CH = 9\text{ cm}$. Tính độ dài đường cao $AH$ (đơn vị: cm). \shortans{$6$}
\end{ex}


\setcounter{ex}{0}
\noindent\textbf{\large PHẦN IV. TỰ LUẬN (3,0 điểm)}

\textit{Thí sinh trình bày chi tiết lời giải vào giấy làm bài.}

\vspace{1mm}
\begin{ex} (1,0 điểm)
Giải các hệ phương trình sau:
\begin{enumerate}[label=\textbf{\alph*)}, leftmargin=0.8cm, itemsep=2pt]
\item (0,5 điểm) $\begin{cases} 2x + 3y = -2 \\ 5x + 2y = 6 \end{cases}$
\item (0,5 điểm) $\begin{cases} 3(x + 7) = 4(y + 3) \\ 4x - 2y + 7 = 0 \end{cases}$
\end{enumerate}
\end{ex}
\vspace{0.5mm}

\begin{ex} (0,75 điểm)
Giải các phương trình sau:
\begin{enumerate}[label=\textbf{\alph*)}, leftmargin=0.8cm, itemsep=2pt]
\item (0,35 điểm) $2(x - 1) - 3 = 5 - 3(x + 2)$
\item (0,4 điểm) $\frac{x + 4}{x - 2} = \frac{1}{x - 2} + 1$
\end{enumerate}
\end{ex}
\vspace{0.5mm}

\begin{ex} (0,75 điểm)
Trong một trận chiến đấu, một máy bay của đối phương bay ở độ cao $1800\text{ m}$. Khẩu pháo cao xạ ngắm chiếc máy bay đó dưới một góc $35^\circ$ so với phương nằm ngang. Tìm khoảng cách từ pháo cao xạ đến máy bay (làm tròn đến mét).
\begin{center}
\begin{tikzpicture}[scale=1.1, line join=round, line cap=round, >=stealth]
  \coordinate (B) at (0,0); % Khau phao
  \coordinate (C) at (4,0); % Hinh chieu may bay tren mat dat
  \coordinate (A) at (4,2.8); % May bay
  \draw[thick, -] (B) -- (C) node[midway, below] {Mặt đất};
  \draw[thick, dashed] (A) -- (C) node[midway, right] {$h = 1800\text{ m}$};
  \draw[thick, color=red!80!black] (B) -- (A) node[midway, above left] {Khoảng cách $d$};
  \draw (3.7,0) -- (3.7,0.3) -- (4,0.3);
  \draw[->, thick, blue] (1,0) arc (0:35:1);
  \node at (1.3,0.3) [blue] {$35^\circ$};
  \filldraw[black] (B) circle (1.5pt) node[below left] {Pháo $(B)$};
  \filldraw[black] (A) circle (1.5pt) node[above] {Máy bay $(A)$};
  \filldraw[black] (C) circle (1.5pt) node[below right] {$(C)$};
\end{tikzpicture}
\end{center}
\end{ex}
\vspace{0.5mm}

\begin{ex} (0,5 điểm)
Cho $x, y$ là hai số thực thỏa mãn điều kiện $x^2 + 2y^2 + 2xy + 7x + 7y + 10 = 0$. Tìm giá trị lớn nhất và giá trị nhỏ nhất của biểu thức $P = \frac{2x + 2y - 3}{x + y + 6}$.
\end{ex}
\vspace{0.5mm}


\enlargethispage{1.5cm}
\begin{center}
\rule{4cm}{0.4pt}\\
\textit{--- HẾT ---}
\end{center}
\end{document}
